\begin{proof}
The sectional calculation follows directly from the
standard analytic form of a $\wt E_7$ simple-elliptic germ,
\[
 z^2+f_4(x,y)=0,
\]
where $f_4$ is a square-free binary quartic
\cite{Laufer1977,Pinkham1977}.  On a general hyperplane $z=ax+by$, the
quadratic term has rank one and the restriction of $f_4$ to its null line is
nonzero.  The analytic splitting lemma therefore reduces the plane-curve
germ to $s^2+t^4=0$, namely type $A_3$, whose Milnor number is $3$.
The Lefschetz formula for the local Euler obstruction
\cite[Theorem~3.1]{BrasseletLeSeade2000} identifies $\operatorname{Eu}_X(0)$
with the Euler characteristic of this generic linear Milnor fibre.  Hence
\[
 \operatorname{Eu}_X(0)=1-\mu(A_3)=1-3=-2.
\]
On the other hand, the stalk of $R\rho_*\QQ_S[2]$ at the origin has
Euler characteristic $e(E_{\exc})=0$, whereas the summand $\delta_0$ has
stalk Euler characteristic $1$.  Thus \eqref{eq:semismall} gives
$\chi((\IC_X)_0)=-1$.  The Dubson--Kashiwara local index formula
\cite{Kashiwara1985} therefore gives
\[
 \operatorname{CC}(\IC_X)=\Lambda_X+m\Lambda_0,
 \qquad -1=-2+m,
\]
hence $m=1$.  The Gauss map of $\Lambda_0=T_0^*P$ is the identity on
$\PP(T_0^*P)$ and has degree one.  The Franecki--Kapranov index theorem
\cite[Theorem~1.3]{FraneckiKapranov2000}, together with
\eqref{eq:chi-seven}, then gives
\[
 7=\deg(\gamma_X)+1.
\]
Thus
\begin{equation}
 \deg(\gamma_X)=6.
\label{eq:gauss-degree}
\end{equation}
\end{proof}

\subsection{The irregularity and the generic-twist Hodge vector}

The exceptional elliptic curve contributes to $H^1(S)$.  We compute this
contribution before considering twists.

\begin{proposition}[Irregularity of the resolution]
\label[proposition]{prop:resolution-irregularity}
The minimal resolution $S$ satisfies
\[
 p_g(S)=5,\qquad q(S)=4.
\]
Consequently
\[
 b_1(S)=8,\qquad b_2(S)=22,\qquad h^{1,1}(S)=12.
\]
\end{proposition}

\begin{proof}
Recall from \eqref{eq:prym-polarization} that
\[
 K(L):=\ker(\varphi_L)=\pi^*E[3]
\]
has order nine.  Every point of $K(L)$ lies on $X$.  For
$\alpha\in E[3]$, the projection formula and
$\pi_*\mathcal O_C=\mathcal O_E\oplus\eta^{-1}\oplus\eta^{-2}$ give
\[
\begin{aligned}
 h^0(C,\kappa\otimes\pi^*\alpha)
 &=h^0(E,\eta\otimes\alpha)+h^0(E,\alpha)
   +h^0(E,\eta^{-1}\otimes\alpha).
\end{aligned}
\]
The first term is one because $\deg\eta=1$; the other two vanish when
$\alpha\ne\mathcal O_E$.  For $\alpha=\mathcal O_E$ the same expression is
two.  Thus
\begin{equation}
 K(L)\subset X.
\label{eq:kernel-on-theta}
\end{equation}

Let $s\in H^0(P,L)$ be a defining section of $X$.  The theta group of $L$
acts on the three-dimensional space $H^0(P,L)$ through its irreducible
Schr\"odinger representation; see, for example,
\cite[Chapter~6]{BirkenhakeLange2004}.  For $a\in K(L)$ choose a lift of
translation by $a$ to the theta group.  The corresponding translate of $s$
vanishes at the origin, because its value there is a nonzero scalar multiple
of $s(a)$ and \eqref{eq:kernel-on-theta} gives $s(a)=0$.  The orbit of the
nonzero vector $s$ spans $H^0(P,L)$ by irreducibility.  Hence
\begin{equation}
 t(0)=0\qquad\text{for every }t\in H^0(P,L).
\label{eq:basepoint-origin}
\end{equation}
In other words, the origin is a base point of the complete linear system
$|L|$.

Adjunction gives $\omega_X\simeq L|_X$, and the divisor sequence is
\[
 0\longrightarrow\mathcal O_P\xrightarrow{\,s\,}L
 \longrightarrow\omega_X\longrightarrow0.
\]
Since $L$ is ample on the abelian threefold, $H^i(P,L)=0$ for $i>0$.
Therefore
\begin{equation}
 0\longrightarrow \CC s\longrightarrow H^0(P,L)
 \longrightarrow H^0(X,\omega_X)
 \longrightarrow H^1(P,\mathcal O_P)\longrightarrow0,
\label{eq:canonical-sections-exact}
\end{equation}
so $h^0(X,\omega_X)=2+3=5$.

All five canonical sections vanish at the singular point.
The two-dimensional subspace coming from $H^0(P,L)/\CC s$ does so by
\eqref{eq:basepoint-origin}.  For the remaining three directions, let $v$ be a translation-invariant
vector field on $P$.  Choose local frames of $L$ with transition functions
$g_{ij}$ and write the defining section as local functions $s_i$, so
$s_i=g_{ij}s_j$.  Differentiating gives
\[
 v(s_i)-g_{ij}v(s_j)=v(g_{ij})s_j.
\]
The right-hand side vanishes on $X$, so the restrictions of the $v(s_i)$
glue to a well-defined section
$\partial_v s\in H^0(X,L|_X)$.  Under the boundary map in
\eqref{eq:canonical-sections-exact}, this section maps, up to the usual sign
convention, to
\[
 \iota_v c_1(L)\in H^1(P,\mathcal O_P),
\]
because the associated \v{C}ech cocycle is $v(g_{ij})/g_{ij}$.  This is
exactly the differential at the origin of the polarization morphism
$\varphi_L:P\to\operatorname{Pic}^0(P)$.  Since $\varphi_L$ is an isogeny,
\[
 H^0(P,T_P)\longrightarrow H^1(P,\mathcal O_P),
 \qquad v\longmapsto\iota_vc_1(L),
\]
is an isomorphism.  The three sections $\partial_v s$ therefore map onto the
last term of \eqref{eq:canonical-sections-exact}.  They also vanish at the
origin because $0$ is a singular point of $X$, so $ds(0)=0$.  Thus every
section of $\omega_X$ vanishes at $0$.

Since $K_X$ is Cartier, write
$K_S=\rho^*K_X+aE_{\exc}$.  Adjunction on the exceptional elliptic curve gives
$(K_S+E_{\exc})\cdot E_{\exc}=0$, while
$\rho^*K_X\cdot E_{\exc}=0$ and $E_{\exc}^2=-2$.  Thus $a=-1$ and
\begin{equation}
 K_S=\rho^*K_X-E_{\exc}.
\label{eq:canonical-resolution}
\end{equation}
The ideal appearing after resolution is also explicit.  Since
$X$ is normal and $\rho$ is a resolution,
$\rho_*\mathcal O_S=\mathcal O_X$.  Push forward
\[
 0\longrightarrow\mathcal O_S(-E_{\exc})\longrightarrow\mathcal O_S
 \longrightarrow\mathcal O_{E_{\exc}}\longrightarrow0.
\]
The exceptional curve is connected and projective, so its global regular
functions are constant.  Since $\rho(E_{\exc})=\{0\}$, the induced map
\[
 \mathcal O_X\longrightarrow\rho_*\mathcal O_{E_{\exc}}
\]
is evaluation at the reduced point $0$ and has kernel $\mathfrak m_0$.
Left exactness of $\rho_*$ gives
\[
 \rho_*\mathcal O_S(-E_{\exc})=\mathfrak m_0.
\]
By the projection formula,
\[
 H^0(S,K_S)=H^0(X,\omega_X\otimes\mathfrak m_0).
\]
The preceding paragraph shows that the right-hand side is all of
$H^0(X,\omega_X)$, and therefore $p_g(S)=5$.  Finally
\cref{lem:surface-invariants} gives $\chi(\mathcal O_S)=2$, so
\[
 q(S)=1+p_g(S)-\chi(\mathcal O_S)=4.
\]
The stated Betti and Hodge numbers now follow from $e(S)=8$.
\end{proof}

Let $j:S\xrightarrow{\rho}X\hookrightarrow P$.  A unitary rank-one character
of $\pi_1(P)$ determines a unitary local system $\mathcal L_\alpha$ and a
holomorphic line bundle $\alpha\in\operatorname{Pic}^0(P)$.  The same symbol $\alpha$ will also denote the latter.

